\section{Capítulo III: Situación Actual, Problemática y Transferencia.}

\subsection{El problema de la subutilización estructural.}
\subsubsection{Inflexibilidad tarifaria frente a la dinámica del mercado privado.}

Según el informe N° 198-GPRC/2017 del Organismo Supervisor de Inversión Privada en Telecomunicaciones 
(OSIPTEL), la fijación de la tarifa de \$27 dólares por Mbps tuvo origen en el planteamiento 
incial del Ministerio de Transportes y Comunicaciones(MTC) bajo el marco de Ley N° 29904.  
Dicha tarifa regulada debía ser pagada directamente a la empresa concesionaria Azteca Comunicaciones Perú
S.A.C. por parte de los operadores que contrataran la capacidad de transporte.
El MTC partió de supuestos de costos y demanda sobredimensionados, lo que colocó a la red dorsal 
en una clara desventaja competitiva.
Según lo expuesto en el panel de DNconsultoresTV (2026), mientras la tarifa estatal 
estaba congelada por contrato, los operadores privados desplegaron su propia infraestructura y bajaron los
precios agresivamente en el mercado, llegando a ofrecer costos de un dígito por Mbps. En consonancia con ello,
Christian Chee argumentó que la RDNFO no contempla descuentos por volumen ni flexibilidad geográfica, por lo que 
las empresas privadas solo la utilizan como último recurso cuando desean extender sus servicios a rutas donde no 
les saldría a cuenta tender sus propias redes de fibra óptica
\subsubsection{El quiebre del modelo de Operador Neutro.}

A pesar del gasto global de más de dos mil millones de dólares, a lo largo del tiempo, en 
este macro despliegue técnico, su modelo financiero como ``operador de operadores`` neutro fracasó.
Al no alcanzar la demanda proyectada por el MTC, el porcentaje operativo real en 2023 según infobae se estancó en 8\%  
de su capacidad total. Además, la brecha fiscal conforme pasa el tiempo se vuelve cada vez más inmanejablen, mientras que el mantenimiento
y funcionamiento de la red le cuesta al estado cerca de setenta millones de dólares, la recaudación
efectiva debido al escaso uso llega a los doce millones de dólares.
Como señala el análisis sectorial de Telesemana(2023), bajo este modelo la ganancia es extremadamente 
baja y los beneficiarios reales son mínimos, 
convirtiendo un activo de alta confiabilidad en un ``elefante blanco`` altamente deficitario.

\subsection{Proceso de caducidad del contrato con Azteca Comunicaciones.}
\subsubsection{Intentos de renegociación y adendas tarifarias.}

El contrato firmado entre el Ministerio de Transportes y Comunicaciones (MTC) y la empresa Azteca 
Comunicaciones se diseñó con una tarifa rígida. Bajo este modelo, la empresa no corría ningún riesgo 
financiero, ya que el Estado le aseguraba por completo el pago de sus costos operativos y de mantenimiento.\\
Este modelo subvencionista generó un fuerte desincentivo comercial, Azteca rechazó las sugerencias del 
Banco Mundial para flexibilizar las tarifas y ofrecer descuentos por volumen. La empresa argumentó que hacer 
esos cambios después de la adjudicación violaría las leyes de las APP (Asociaciones Público-Privadas) y las 
reglas de la licitación. Según Azteca, cambiar las reglas con las que ganaron el contrato habría sido injusto 
y legalmente riesgoso, ya que rompería la igualdad de condiciones frente a las demás empresas. \\
Para intentar salvar el proyecto y frenar el derroche fiscal, más adelante se propuso unificar 
la red dorsal con la red de transporte. Sin embargo, esta idea no tenía sustento, ya que no había estudios 
ni pruebas sólidas de que la red de transporte pudiera sostenerse sola. De haberse aplicado, la fusión habría 
empeorado los problemas de ambas redes: integrar una infraestructura tan compleja solo habría creado una «megadorsal» 
difícil de manejar y de menor calidad.